\chapter{地球化学}
\section{背景：水を介して岩石・大気・生物を結ぶ}
\label{theory:geochemistry:boundary}
地球化学は、元素がどこに存在し、どの形態で移り、何が濃度を変えるかを扱う。
物理専攻者には、保存則と輸送方程式が入口になる。ただし元素の総量を保存しても、
溶存イオン、吸着、沈殿、気体、生体の間で配分が変わる。濃度の変化を水量の変化だけで説明できない。
降水が土壌を通り、岩石と反応し、地下水となり、河川から海へ運ぶ流れに沿って、
酸塩基・電子授受・溶解平衡・反応速度を結び付ける。

この初稿は地球化学の登録22問を原典19頁画像で確認した範囲である。
各問の公式正答と独立解答、各肢の成立条件を分け、原典の不十分な名称や強い一般化を改変せず記録する。
同位体の詳細は同位体地球化学章、地下水の水理は地下水章にも接続する。

\section{貯留・流量・濃度：何を比較しているか}
\label{theory:geochemistry:reservoir}
ある区画の水量を \(V\)、溶存成分濃度を \(C\) とすれば質量は \(M=VC\) である。
水の流量 \(Q\) は体積/時間、溶質の移動率 \(QC\) は質量/時間で、貯留量とは異なる。
区画を十分混合と近似すると
\[
 \frac{d(VC)}{dt}=\sum_j Q_jC_j-Q_{\rm out}C+P-L
\]
である。\(P,L\) は反応・供給・除去の質量/時間である。
蒸発は水を除いて塩を残すため、成分を追加しなくても \(C\) が上がる。
沈殿や吸着は溶存量を下げるが、系内の元素全量を必ず減らすわけではない。

USGSの代表貯留表は総水量約 \(1.386\times10^9\,\mathrm{km^3}\)、海約96.54\%、
河川約2120、大気約12900、全地下水約2340万 \(\mathrm{km^3}\) とする\cite{gc_water}。
大きな年間河川流出量を、河川内の瞬時貯留へ代入しない。
地下水は提示された湖・河・湿地・土壌より大きいが、取り出せる淡水資源の量と全地下水量は別である。
\begin{tcolorbox}[title=「約70\%」の分母]
氷河・雪氷は淡水貯留の約68.7\%である。海洋以外の全水に塩性地下水まで含めると、
同じ約70\%にはならない。平成26年度問25と令和元年度再試験問18は「陸水」という語を使うため、
教材では原語を残して淡水を分母とした意図と範囲疑義を併記する。
\end{tcolorbox}
平均滞留時間 \(V/Q\) は定常・同じ区画の出入りを近似する指標であり、
個々の水分子の年齢が全てその値になるという意味ではない。

\section{汚濁指標と起源を追うトレーサー}
\label{theory:geochemistry:indicators}
BODは生物分解に伴う酸素要求量で、分解可能な有機負荷を追う。
SSは浮遊粒子、全窒素は窒素成分の合計、大腸菌群数は出題当時の衛生汚濁の指標である。
これらは同じ「水質」でも対象が違う。BODを全有機炭素量に、SSを溶存量に置き換えない。
測定対象・方法・単位を確かめる必要がある\cite{gc_hem}。
水素同位体比は蒸発、降水、混合等で変わり、水の起源や移動を追う。
比が高いだけで汚濁量が多いとは判断できない。
汚染研究に同位体を使うことと、同位体比が一般的な汚濁負荷の直接指標であることは別である。

\section{河川水は海水を薄めたものだけではない}
\label{theory:geochemistry:river}
雨に取り込まれた海塩、土壌の反応、岩石風化、温泉・鉱山・人為排水、蒸発、海水混合が組成を変える。
世界平均河川の主要イオンにはCaと炭酸水素が重要で、NaClが世界の大半を支配するという一般化は誤る。
沿岸の降水に海塩が多いことを、全河川の溶存成分の起源へ拡張しない。

北野・加藤の1970年論文の比較は、小林1960とLivingstone1963の集計を用いる。
表4では日本の蒸発残渣74.8mg/L、溶存SiO2の割合26.84\%、世界100mg/L、11.67\%である。
割合を濃度へ戻すと
\[
 C_{\mathrm{SiO_2,Japan}}=74.8\times0.2684\simeq20.1\,\mathrm{mg/L},\qquad
 C_{\mathrm{SiO_2,world}}=100\times0.1167=11.67\,\mathrm{mg/L}
\]
となる。同論文は多雨・蒸発濃縮の小ささと火山地質を説明する\cite{gc_river1970}。
古い試料群の一般傾向を、現在の全流域が必ず同じ順位であるという命題にしない。
Si表示とSiO2換算表示も分ける。

\section{炭酸平衡：石灰岩を溶かす条件と析出させる条件}
\label{theory:geochemistry:carbonate}
土壌呼吸で供給されたCO2が水に溶けると炭酸系を作る。
\[
 \mathrm{CO_2(g)\rightleftharpoons CO_2(aq)},\qquad
 \mathrm{CO_2(aq)+H_2O\rightleftharpoons H^++HCO_3^-},\qquad
 \mathrm{HCO_3^-\rightleftharpoons H^++CO_3^{2-}} .
\]
第一式の気液分配と後二式の酸解離は異なる平衡である。
ここで第二式のCO2には水和した炭酸を合わせた慣用的な表現を用いた。
各平衡定数は温度・活量基準に依存する。
石灰岩の代表鉱物である方解石について
\[
 \mathrm{CaCO_3+CO_2+H_2O\rightleftharpoons Ca^{2+}+2HCO_3^-}
\]
が、Caと炭酸水素を増やす反応を示す。
CO2を多く含む水は溶解側へ、洞窟でCO2を失う水は析出側へ動き、
鍾乳石や石筍を形成する\cite{gc_hem,gc_cave}。
CO2添加を普通の洞窟での析出原因とする令和6年度問22の肢2は方向が逆である。
\begin{tcolorbox}[title=原典の水質名を科学的な文へ書き換えない]
平成25年度問31と令和元年度問20は「炭酸カルシウム型」と呼ぶ。
通常、主要溶存種のCaとHCO3に基づく名称は炭酸水素カルシウム型、すなわちCa-HCO3型である。
固体CaCO3と、Ca・HCO3が別々に溶存する水を混同しない\cite{gc_hem}。
原典はそのまま保存し、名称の疑義を問題ごとに保持する。
\end{tcolorbox}

\section{採水後に変わるpH：測った時点も条件である}
\label{theory:geochemistry:sampling}
地下水が大気より高いCO2分圧の土壌と接した場合、採取・攪拌で脱気するとpHが上昇し得る。
したがって現場の非攪拌水と室内放置水は同じ組成の測定とは限らない。
ただし酸化反応、沈殿、空気からのCO2吸収等もあり、全試料で必ず上がるという規則ではない。
温度、曝気、保存時間を記録して現場条件を優先する\cite{gc_hem}。
pHは水素イオン活量の負の常用対数であり、濃度との同一視は希薄な近似である。
炭酸系の平衡は酸化還元を伴わなくてもpHを変える。

\section{活量：濃度だけで平衡を記述できない理由}
\label{theory:geochemistry:activity}
理想溶液では化学ポテンシャルの組成依存を濃度比で表せるが、イオンは互いの電場に影響される。
規定した標準状態に対し
\[
 \mu_i=\mu_i^\circ+RT\ln a_i,\qquad a_i=\gamma_i\frac{c_i}{c^\circ},\qquad
 I=\frac12\sum_i c_i z_i^2
\]
と置く。\(a_i,\gamma_i\) は無次元、\(z_i\) は価数で、イオン強度 \(I\) の濃度基準を揃える。
二価イオンは同じ濃度でも \(z_i^2=4\) の重みを持つ。
Debye--Huckelの極限則は希薄溶液の静電効果から
\[
 \log_{10}\gamma_i\simeq-Az_i^2\sqrt I
\]
という形を与える。\(A\) の値と単位は温度・溶媒・濃度基準に依存する。
濃厚塩水ではこの極限形を外挿せず、より適切な活量モデルを使う。
単一イオン活量の分離には慣例が必要である\cite{gc_debye}。
これは鉱物の風化抵抗順位そのもののモデルではない。

\section{平衡でどこまで溶けるかと、どれほど速く溶けるか}
\label{theory:geochemistry:equilibrium}
固定した温度・圧力・化学条件で、残存する固相と水の化学ポテンシャルが釣り合うと溶解平衡となる。
方解石ならイオン活動度積 \(a_{\mathrm{Ca^{2+}}}a_{\mathrm{CO_3^{2-}}}\) を溶解度積と比較する。
溶存総炭素とCO3の活量はpHによる種分布のため同じではない。
少量の固相が全部溶けた場合、観測濃度が飽和濃度とは限らない。
開放流動、別相生成、過飽和の維持もあり、「時間がたてば必ず飽和平衡」とはしない。
平衡は到達先を示すが、表面積、温度、反応経路、飽和からの距離が速度を変える\cite{gc_hem}。
令和4年度問20では、活量の理論・平衡・経験的抵抗性・速度の対象を区別する。

\section{地殻の存在度：対象と質量基準を固定する}
\label{theory:geochemistry:crust}
地殻は酸化物・ケイ酸塩を多く含み、質量基準ではO、Si、Alが上位である。
USGSの代表表はO46.60、Si27.72、Al8.13、K2.59、Ti0.44、Mn0.100、Cl0.020\%を示す\cite{gc_crust}。
平成25年度問10のCl>Kは逆である。
各小数は地殻モデルや試料集合に依存し、上部大陸地殻、全地殻、地球全体、海水では組成が違う。
化学式中の原子比と岩石中の質量存在度も別の比較である。

\section{酸化還元：電子供与体と受容体を先に置く}
\label{theory:geochemistry:redox}
電子を与える物質は自身が酸化されるので還元剤、受け取る物質は自身が還元されるので酸化剤である。
酸素供給の乏しい地下水では、利用可能な反応物と微生物活動に応じ、
酸素、硝酸、Mn・Fe酸化物、硫酸、CO2等が電子受容体として使われる\cite{gc_redox,gc_redox2009}。
これは各反応が常にきれいな順で完了するという保証ではない。
Fe(III)の固相が還元されるとFe(II)が溶存し得るため、還元地下水はFeMnに富みやすい。
湧出して酸素を得るとFe(III)水酸化物等が茶褐色に析出する。
沈殿の酸化状態を、その水が地下で持っていた状態と混同しない。

基準電極に対して測る電位を揃える。標準水素電極を0とする \(E_h\) と、
測定器の参照電極に対する表示は換算が必要である。
酸素の還元半反応
\[
 \mathrm{O_2+4H^++4e^-\rightleftharpoons2H_2O}
\]
にNernst式を用い、水活量1、気体の分圧を標準圧で無次元化すると、25℃では
\[
 E_h\simeq1.229-0.05916\,\mathrm{pH}
       +0.01479\log_{10}p_{\mathrm{O_2}}\quad(\mathrm V)
\]
である。pH7、空気 \(p_{\mathrm{O_2}}=0.21\) で約+0.805Vとなり、-0.250Vとは符号が異なる。
これは理想的な酸素/水の境界であり、白金電極の実測値が必ず一致するという説明ではない。
複数反応の混合電位、反応の遅さ、電極応答がある\cite{gc_hem}。
地下の長い滞留でも、酸化可能な物質との反応なしに酸素が消えるわけではない。
火山温泉の還元性と、湧出後の空気混合も区別する。

\section{硫化物の酸化と酸性鉱山排水}
\label{theory:geochemistry:mine}
代表例として黄鉄鉱の酸化は
\[
 \mathrm{FeS_2+\tfrac72O_2+H_2O\longrightarrow Fe^{2+}+2SO_4^{2-}+2H^+}
\]
と書ける。硫酸イオンと酸を生じ、さらに鉄の酸化・加水分解等が関わる。
露出面積、酸素供給、微生物、炭酸塩等の中和能が結果を変える。
硫化物があるだけで全ての排水が同じpHになるとは限らず、
中性や塩性で溶存金属の多い排水もある\cite{gc_mine,gc_mine2015}。
Ca-SO4型を生じるには、硫酸だけでなくCaの供給も必要である。

\section{化学的風化：一次鉱物から二次鉱物へ}
\label{theory:geochemistry:weathering}
地表の水・温度・酸素条件は、深部やマグマから形成された鉱物の条件と違う。
通常の造岩鉱物では、かんらん石、輝石、黒雲母、石英の順に化学的風化への抵抗が増す。
これは環境と鉱物の反応をまとめた経験系列で、物理的破砕の順位や全試料の反応速度を保証しない。
方解石は炭酸を含む水で比較的速く溶けるが、飽和に近い水では速度が低下する。

塩基性陽イオンと一部のSiが溶脱すると、残るAl・Feが酸化物・オキシ水酸化物や粘土を作る。
強い溶脱が進む湿潤熱帯ではカオリナイトが生じやすい。
スメクタイト等との違いは層状構造・組成と溶液条件の違いである。
岩種、水岩比、温度、生物、期間が重要で、気候名だけで粘土を決めない\cite{gc_clays,gc_hem}。
酸化的な通常土壌でのAl・Fe(III)沈殿と、酸性水・錯体・還元Fe(II)の溶存を区別する。

\section{炭素循環：速い生物過程と遅い岩石過程}
\label{theory:geochemistry:carbon}
光合成でCO2を有機炭素に取り込み、呼吸・分解で戻す速い循環と、
風化・炭酸塩堆積・有機物埋没・変成・火山による遅い循環を分ける\cite{gc_carbon}。
貯留が大きいことは、現在の吸収速度も大きいこととは別である。
海洋の溶存無機炭素等の貯留は大気よりはるかに多く、半分や同程度という関係ではない。

単純化したCaケイ酸塩の風化は
\[
 \mathrm{CaSiO_3+2CO_2+H_2O\longrightarrow Ca^{2+}+2HCO_3^-+SiO_2}
\]
である。海でCaCO3が沈殿してCO2の一部が戻る反応を加えると、
\[
 \mathrm{CaSiO_3+CO_2\longrightarrow CaCO_3+SiO_2}
\]
という正味の固定を表す。岩石を単一鉱物とみなした収支例で、実際の速度をこの式だけで求めない。
土壌の高いCO2は炭酸を供給し風化を促進し得る。
初期のCO2濃度の低下に風化と生物活動が関与したという地史的な説明は、
今日の大気が既に完全な定常状態にあるという主張ではない。

自然の温室効果に大きく寄与する水蒸気と、人為的な追加放射強制力のCO2を区別する。
CO2の主要な温室作用は赤外の吸収・放射であり、紫外の吸収との取り違えを避ける。
メタン生成は通常嫌気的な代謝で、光合成ではない。例として
\[
 \mathrm{CO_2+4H_2\longrightarrow CH_4+2H_2O}
\]
は還元条件での基質と電子収支を示す。原典の「メタン生成菌」という呼び方を、
生物の分類や全生成経路の説明にそのまま拡張しない。
\begin{tcolorbox}[title=炭素の収支を閉じるために]
平成23年度問19と平成28年度問27の正答肢は、
風化の除去作用が光合成・分解・呼吸で長期に均衡すると書く。
重要な過程を挙げる意図は理解できるが、遅い循環の火山・堆積・埋没等を省く。
教材では原語と公式を保存し、定常仮定の範囲を疑義として残す。
\end{tcolorbox}

\section{窒素循環：窒素の反応と生物の炭素源は別軸である}
\label{theory:geochemistry:nitrogen}
窒素固定はN2から反応性窒素を作り、硝化は還元態窒素を酸化し、脱窒は硝酸等を還元して気体へ戻す。
硝化の代表的な全反応は
\[
 \mathrm{NH_4^++2O_2\longrightarrow NO_3^-+2H^++H_2O}
\]
である。アンモニアから亜硝酸、さらに硝酸の段階を分ける。
固定ではまずN2の還元によりアンモニアを得てアミノ酸へ同化する。
雷や燃焼の熱的NOxは非生物固定だが、燃料中に既に存在する窒素の再排出と新規N2固定は違う。

有機物を電子供与体とする脱窒の収支例を \(\mathrm{CH_2O}\) で単純化すると
\[
 \mathrm{5CH_2O+4NO_3^-+4H^+\longrightarrow2N_2+5CO_2+7H_2O}
\]
となる。これを全ての脱窒生物の炭素源へ一般化しない。
無機硫黄を電子供与体とする例は
\[
 \mathrm{5S+6NO_3^-+2H_2O\longrightarrow3N_2+5SO_4^{2-}+4H^+}
\]
で、硫黄酸化と硝酸還元がエネルギーを供給する。
1989年の一次実験は、HCO3を無機炭素源、還元硫黄を電子供与体とする嫌気的な自養脱窒菌の集積を示す\cite{gc_sulfur_denit}。
電子供与体、受容体、同化する炭素源を別々に答えれば、硝化と脱窒の二語だけで栄養型を決める誤りを避けられる。
\begin{tcolorbox}[title=令和2年度問19の公式と科学的例外]
公式は肢5を不適切とする。独立初回も5であった。
ただし化学合成と脱窒を組み合わせて無機炭素を同化する過程は実在する。
硝化と従属栄養脱窒の対比という設問意図を説明し、肢5を全生物に一律で否定しない。
疑義と査読待ちを保持する。
\end{tcolorbox}

\section{溶質輸送：移流・拡散・分散を分ける}
\label{theory:geochemistry:transport}
同じ水とともに移動する移流、熱運動による分子拡散、流れの速度と経路の差による機械分散を分ける。
濃度 \(C\)、平均間隙流速 \(v\)、一定の拡散分散係数 \(D\) について、
一様な一次元、吸着と反応なしという模型では
\[
 J=vC-D\frac{\partial C}{\partial x},\qquad
 \frac{\partial C}{\partial t}+v\frac{\partial C}{\partial x}
 =D\frac{\partial^2 C}{\partial x^2}
\]
が流束と連続の式から得られる。
Fick則の比例する量は濃度勾配に対する流束であり、個々の分子の一定の流速ではない。
点状の初期分布を無限領域へ流せば中心は \(vt\)、分散は \(2Dt\) となり、移動と広がりを区別できる。

縦方向には \(D_L=\alpha_L|v|+D_m^\ast\)、横方向にも対応する式を使う。
\(\alpha_L,\alpha_T\) は長さ、\(D_m^\ast\) は媒体内の有効分子拡散係数で面積/時間である。
通常、横分散長は縦より小さい。野外では室内より大きい不均質性を含むため、
実験室の値をそのまま流域へ外挿できない。
59現場を比較した1992年研究は、観測尺度と分散長の傾向に加え、
同じ尺度でも大きな散らばりと信頼度差があるため普遍直線が不適切だとする\cite{gc_dispersion}。

\section{電気伝導率：イオンの量と動きやすさを測る}
\label{theory:geochemistry:conductivity}
水の電気伝導はイオンの移動による。希薄な独立寄与の近似なら
\[
 \kappa\simeq\sum_i\lambda_i c_i,\qquad
 \kappa=\frac1\rho=\frac{\ell}{RA}
\]
である。\(\lambda_i\) はイオン種に応じたモル伝導寄与、
\(\rho\) は抵抗率、\(R\) は長さ \(\ell\)・断面積 \(A\) の試料の電気抵抗である。
濃度の単位と \(\lambda_i\) の単位を合わせる。
1cm立方体では数値として \(1/R\) に対応するが、伝導率の単位はS/cmで、伝導度のSと異なる。
SIのS/mへの換算は \(1\,\mathrm{S/cm}=100\,\mathrm{S/m}\) である。

同じ天然水の通常温度域では温度上昇でイオン移動度が増し、伝導率は上がる。
比較には温度を揃え、25℃へ補正した値を用いる\cite{gc_hem,gc_conductance,gc_conductance1994}。
海水、河川水、雨の順位は塩類量の一般傾向で、例外はある。
同じ組成比なら濃度と蒸発残留物に関係するが、組成が違えば単一係数で総溶存量へ換算できない。

\section{水の同位体：正の相関と履歴の識別}
\label{theory:geochemistry:isotope}
標準に対する同位体比 \(R\) の偏差を
\[
 \delta=\left(\frac R{R_{\rm std}}-1\right)1000\quad(\text{千分率})
\]
で表す。重い同位体は蒸発・凝結の際に軽いものと同じ比率で移らない。
HとOの分別が連動するため、代表的な全球天水線は
\[
 \delta^2\mathrm H\simeq8\delta^{18}\mathrm O+10
\]
という正の相関である\cite{gc_meteoric}。
絶対同位体比が厳密な反比例であるという意味ではない。
局地天水線の傾き・切片は違い、蒸発や混合は天水線から逸脱させる。
逐次凝結で重い同位体が除かれると、高緯度へ向かう水蒸気と降水は重い同位体に乏しくなる一般傾向がある。
起源識別と単独の汚濁量評価を区別する。

\section{生物によるケイ酸の取り込みと、殻の溶解}
\label{theory:geochemistry:silica}
珪藻は溶存ケイ酸を取り込み、殻を形成する。
成長時の取り込みを、河川ケイ酸濃度を上げる直接作用として説明するのは逆である\cite{gc_hem,gc_diatom}。
しかし死後の殻の溶解はケイ酸を水へ戻す。
水中濃度の変化は、風化供給、取り込み、溶解、輸送、沈殿等の差で決まり、個々の過程を区別する。

海洋の粘土生成である逆風化は、アルカリ度とカチオンを消費してCO2を生じ、
pHを下げ得る。これは陸上の単純なケイ酸塩風化の酸消費とは違う。
2025年の原著研究はこの反応と生物起源シリカの役割を検討する\cite{gc_reverse}。
後年の科学説明であって、平成27年原典の意図を直接証明したものとはしない。
同問の「ケイ藻によってケイ素が放出されることがある」は生活史の範囲を省くため、
公式5を残しつつ初回4|5と条件疑義を保持する。

\section{自然ヒ素：化学形態・濃度・基準の時点}
\label{theory:geochemistry:arsenic}
ヒ素は鉱物中に存在し、酸化還元、溶解、吸着・脱着、二次鉱物生成により地下水等へ移り得る。
元素の全量、無機のIII・Vの形態、単体、溶存種、鉱物を分ける\cite{gc_hem,gc_redox}。
物理・化学的な取り込みの違いを、単体が一律に無害という説明にしない。
NIOSHの金属ヒ素を含む一次有害性資料は、原典の「毒性はない」という断言を支持しない\cite{gc_arsenic_hazard}。

\begin{tcolorbox}[title=平成30年度問28：単位の誤りと他肢の疑義を分ける]
原典肢2は0.01g/Lと書く。2004施行の水道基準表と現在の表はヒ素量0.01mg/Lである\cite{gc_arsenic2004,gc_arsenic_current}。
0.01g/Lは10mg/Lで1000倍違うため、公式2の根拠は確認できる。
しかし肢3の無毒という断言と肢4の岩種比較・揮発性の因果は別に検証を要する。
原典単位を正しいmg/Lへ校正して誤りを消したり、公式と一致しただけで全肢を確認済みにしたりしない。
\end{tcolorbox}
基準は用途・対象・施行時点に依存する。現在の表の項目数や他成分の変更を2018年へ遡及しない。

\section{CO2回収・貯留：分離と封じ込めの役割}
\label{theory:geochemistry:ccs}
CCSは排出等からCO2を回収し、移送して地中へ貯留する工程である。
アミン吸収や膜分離は回収、深部塩水帯水層等は貯留場所、
CO2-EORは利用と貯留へ接続する\cite{gc_ccs,gc_capture,gc_mea}。
帯水層という一候補だけで地下水流動の問題へ主分類を限定しない。
EORにはCO2以外の方式もあり、界面活性剤の油回収への利用が全てCCSそのものになるわけではない。
平成26年度問18は「関係性が最も薄い」という相対判断で、界面活性剤を無関係と断言する問いではない。
圧力、流動、鉱物反応、封じ込めと漏えいを別の評価課題として扱う。

\section{今回の出題形式と次の確認}
22問では、陸水の組成5問、炭素循環4問、水貯留3問が反復し、
指標、地殻存在度、CCS、窒素、輸送、伝導率、酸化還元、水鉱物、土壌風化、ヒ素が各1問である。
これは今回実際に読んだ問題の集計で、未処理の他分類や将来の出題確率を含む最終頻度ではない。
物理専攻者は、まず質量保存と流束、次に化学種を決める酸塩基・電子授受・活量、
最後に地質・生物・測定時点の条件を加えると、用語の暗記だけに依存せず判断できる。
正答位置が入れ替わった同型でも、方向・単位・分母・限定語を原典で再確認する。
疑義11問はAstra査読対象として保持し、以下の個別解説から条件を追跡できる。
\section{炭素貯留・生物循環と岩石風化}

\subsection{平成23年度 問19}
\label{q:h23-reg-applied-19}

\noindent\textbf{原典：} sources/H23/17\_applied\_science.pdf、PDF物理7頁。

\textbf{公式正答：} 3（single）。

\par\textbf{検討未確定：} 正答の風化・生物循環の重要性は支持されるが、「光合成、有機物分解、呼吸によって長期にバランス」という原典の収支範囲は不十分。火山・堆積・埋没も必要で、人為擾乱下の現在大気は定常でない。原典を修正せず部分校正。

\paragraph{問題の要旨} 炭素の貯留量、自然温室効果、生物過程と地史的風化を比較する。

\Needspace{4\baselineskip}\paragraph{確認した条件・記述（独自要約）}\mbox{}\par

\noindent\textbf{方向：} 正しいもの。\par

\noindent\textbf{対象：} 初期大気・現在の自然温室効果・炭素貯留と循環。\par

\paragraph{問われる核心} 炭素の速い生物循環と遅い岩石循環

\paragraph{必要知識} 炭素の速い生物循環と遅い岩石循環。設問が意図する3は風化と生物活動の重要性であり、初回・公式とも3である。他の候補は初期濃度、水蒸気の役割、メタン生成の酸化還元条件、海洋の大きな貯留量に反する。ただし風化の収支を生物3作用だけで閉じない。短い生物循環と火山・堆積を含む遅い循環を分け、現在の炭素増加を隠す定常仮定にも用いない。

\paragraph{考え方・正答を選ぶ理由（査読待ち初稿）} 設問が意図する3は風化と生物活動の重要性であり、初回・公式とも3である。他の候補は初期濃度、水蒸気の役割、メタン生成の酸化還元条件、海洋の大きな貯留量に反する。ただし風化の収支を生物3作用だけで閉じない。短い生物循環と火山・堆積を含む遅い循環を分け、現在の炭素増加を隠す定常仮定にも用いない。

\Needspace{4\baselineskip}\paragraph{選択肢の検討（査読待ち初稿）}\mbox{}\par

\noindent\textbf{1：不適切} 初期CO2は現在より低い。。初期大気のCO2は現在より高かったとする地史的な理解と逆である。初期という語の全時代の精密濃度を固定しない。\par

\noindent\textbf{2：不適切} 最大の温室効果はCO2、次いで水蒸気。。自然の温室効果では水蒸気の寄与が大きい。人為的な追加放射強制力のランキングを現在大気の全温室効果と混ぜている。\par

\noindent\textbf{3：意図上適切・収支範囲疑義} 風化によるCO2除去は重要で、光合成・分解・呼吸で長期に均衡。。風化による炭素固定と生物の吸収・呼吸・分解は重要である。ただし遅い循環の収支は火山・堆積・有機炭素埋没等も含み、3作用だけで任意時点の均衡が保証されるわけではない。\par

\noindent\textbf{4：不適切} 酸化環境でメタン生成。。メタン生成微生物は通常嫌気的な条件で活動する。酸化的な条件でのメタン酸化と生成を区別する。\par

\noindent\textbf{5：不適切} 大気炭素は海洋炭素の半分。。海洋の溶存無機炭素等は大気炭素よりはるかに大きい。大気が海洋の半分という量関係ではない。\par

\paragraph{再利用する考え方} 物質の貯留量・輸送速度・化学形態を分け、系の境界と反応の前提を確かめる。炭素の速い生物循環と遅い岩石循環を収支・平衡・反応速度のどのモデルで扱うか選ぶ。

\paragraph{共通理論} \hyperref[theory:geochemistry:carbon]{第\ref*{theory:geochemistry:carbon}節}, \hyperref[theory:geochemistry:weathering]{第\ref*{theory:geochemistry:weathering}節}

\par 参考：\cite{gc_carbon,atm_noaa_ghg}。

\paragraph{関連問題} \hyperref[q:h26-reg-applied-27]{H26-REG-APPLIED-27}, \hyperref[q:h28-reg-applied-27]{H28-REG-APPLIED-27}

\par 難易度 3（校準初稿）、優先度未評価。

\subsection{平成26年度 問27}
\label{q:h26-reg-applied-27}

\noindent\textbf{原典：} sources/H26/17\_applied\_science.pdf、PDF物理11頁。

\textbf{公式正答：} 4（single）。

\paragraph{問題の要旨} CO2吸収帯、風化、海洋貯留、地史的変化とメタン生成を比較する。

\Needspace{4\baselineskip}\paragraph{確認した条件・記述（独自要約）}\mbox{}\par

\noindent\textbf{方向：} 最も適切なもの。\par

\paragraph{問われる核心} 炭素循環の貯留・風化・微生物代謝

\paragraph{必要知識} 炭素循環の貯留・風化・微生物代謝。紫外を赤外、減速を促進、同程度を海洋の大きな貯留、光合成を嫌気代謝へそれぞれ区別する。地史的な風化と光合成の作用を述べる4が最も適切である。生物が吸ったCO2が全量恒久貯留するわけではない。

\paragraph{考え方・正答を選ぶ理由（査読待ち初稿）} 紫外を赤外、減速を促進、同程度を海洋の大きな貯留、光合成を嫌気代謝へそれぞれ区別する。地史的な風化と光合成の作用を述べる4が最も適切である。生物が吸ったCO2が全量恒久貯留するわけではない。

\Needspace{4\baselineskip}\paragraph{選択肢の検討（査読待ち初稿）}\mbox{}\par

\noindent\textbf{1：不適切} CO2は紫外に幅広い吸収帯。。主要な温室作用は赤外領域の吸収・放射である。紫外との取り違え。\par

\noindent\textbf{2：不適切} 土壌CO2が高いので風化のCO2除去が減速。。土壌呼吸によるCO2供給は炭酸を作り化学的風化を促す。実際の速度は水流、表面積、温度、飽和度等にも依存する。\par

\noindent\textbf{3：不適切} 大気炭素全量は海洋と同程度。。海洋の炭素貯留は大気よりはるかに多く、同程度ではない。\par

\noindent\textbf{4：適切} 初期CO2は高く、減少へ風化と光合成が作用。。地史的なCO2減少に風化と生物活動が関与する。光合成固定も呼吸で戻る部分と埋没する部分を区別する。\par

\noindent\textbf{5：不適切} 嫌気環境のメタン生成は光合成。。嫌気条件は整合するがメタン生成は光合成ではない。水素とCO2、酢酸等の基質を使う代謝である。\par

\paragraph{再利用する考え方} 物質の貯留量・輸送速度・化学形態を分け、系の境界と反応の前提を確かめる。炭素循環の貯留・風化・微生物代謝を収支・平衡・反応速度のどのモデルで扱うか選ぶ。

\paragraph{共通理論} \hyperref[theory:geochemistry:carbon]{第\ref*{theory:geochemistry:carbon}節}, \hyperref[theory:geochemistry:weathering]{第\ref*{theory:geochemistry:weathering}節}, \hyperref[theory:geochemistry:redox]{第\ref*{theory:geochemistry:redox}節}

\par 参考：\cite{gc_carbon,atm_noaa_ghg,gc_redox}。

\paragraph{関連問題} \hyperref[q:h23-reg-applied-19]{H23-REG-APPLIED-19}, \hyperref[q:r03-reg-applied-21]{R03-REG-APPLIED-21}

\par 難易度 3（校準初稿）、優先度未評価。

\subsection{平成28年度 問27}
\label{q:h28-reg-applied-27}

\noindent\textbf{原典：} sources/H28/17\_applied\_science.pdf、PDF物理11頁。

\textbf{公式正答：} 4（single）。

\par\textbf{検討未確定：} 正答の風化・生物循環の重要性は支持されるが、「光合成、有機物分解、呼吸によって長期にバランス」という原典の収支範囲は不十分。火山・堆積・埋没も必要で、人為擾乱下の現在大気は定常でない。原典を修正せず部分校正。

\paragraph{問題の要旨} 自然温室効果、初期大気、メタン条件、風化と炭素貯留を判定する。

\Needspace{4\baselineskip}\paragraph{確認した条件・記述（独自要約）}\mbox{}\par

\noindent\textbf{方向：} 最も適切なもの。\par

\paragraph{問われる核心} 速い生物循環と遅い炭素循環の収支範囲

\paragraph{必要知識} 速い生物循環と遅い炭素循環の収支範囲。初回・公式とも4であり、地史と生物過程を扱う意図で選ぶ。ただし遅い炭素循環は火山・堆積・埋没も含み、生物3作用だけで風化流束まで均衡するという説明にはしない。H23問19と同じ文言群だが正答位置と問いの方向を原典で再照合した。

\paragraph{考え方・正答を選ぶ理由（査読待ち初稿）} 初回・公式とも4であり、地史と生物過程を扱う意図で選ぶ。ただし遅い炭素循環は火山・堆積・埋没も含み、生物3作用だけで風化流束まで均衡するという説明にはしない。H23問19と同じ文言群だが正答位置と問いの方向を原典で再照合した。

\Needspace{4\baselineskip}\paragraph{選択肢の検討（査読待ち初稿）}\mbox{}\par

\noindent\textbf{1：不適切} 初期CO2は現在より低い。。初期大気のCO2は現在より高かったとする地史的な理解と逆である。初期という語の全時代の精密濃度を固定しない。\par

\noindent\textbf{2：不適切} 最大の温室効果はCO2で水蒸気が次。。自然の温室効果では水蒸気の寄与が大きい。人為的な追加放射強制力のランキングを現在大気の全温室効果と混ぜている。\par

\noindent\textbf{3：不適切} メタンは酸化環境で生成。。メタン生成微生物は通常嫌気的な条件で活動する。酸化的な条件でのメタン酸化と生成を区別する。\par

\noindent\textbf{4：意図上適切・収支範囲疑義} 風化によるCO2除去は重要で、光合成・分解・呼吸で長期に均衡。。風化による炭素固定と生物の吸収・呼吸・分解は重要である。ただし遅い循環の収支は火山・堆積・有機炭素埋没等も含み、3作用だけで任意時点の均衡が保証されるわけではない。\par

\noindent\textbf{5：不適切} 大気Cは海洋Cの半分。。海洋の溶存無機炭素等は大気炭素よりはるかに大きい。大気が海洋の半分という量関係ではない。\par

\paragraph{再利用する考え方} 物質の貯留量・輸送速度・化学形態を分け、系の境界と反応の前提を確かめる。速い生物循環と遅い炭素循環の収支範囲を収支・平衡・反応速度のどのモデルで扱うか選ぶ。

\paragraph{共通理論} \hyperref[theory:geochemistry:carbon]{第\ref*{theory:geochemistry:carbon}節}, \hyperref[theory:geochemistry:weathering]{第\ref*{theory:geochemistry:weathering}節}

\par 参考：\cite{gc_carbon,atm_noaa_ghg}。

\paragraph{関連問題} \hyperref[q:h23-reg-applied-19]{H23-REG-APPLIED-19}

\par 難易度 3（校準初稿）、優先度未評価。

\subsection{令和3年度 問21}
\label{q:r03-reg-applied-21}

\noindent\textbf{原典：} sources/R03/17\_applied\_science.pdf、PDF物理8頁。

\textbf{公式正答：} 2（single）。

\paragraph{問題の要旨} 炭素貯留、地史的風化、生物代謝とCO2の吸収帯を比較する。

\Needspace{4\baselineskip}\paragraph{確認した条件・記述（独自要約）}\mbox{}\par

\noindent\textbf{方向：} 最も適切。\par

\paragraph{問われる核心} 貯留・反応・放射を区別する炭素循環

\paragraph{必要知識} 貯留・反応・放射を区別する炭素循環。地史的な濃度減少への複数の関与を述べる2が適切。海洋と大気を同量にせず、メタン生成を光合成にせず、紫外と赤外を取り違えず、土壌CO2による炭酸供給の風化促進を理解する。風化速度・炭素埋没の時間尺度も意識する。

\paragraph{考え方・正答を選ぶ理由（査読待ち初稿）} 地史的な濃度減少への複数の関与を述べる2が適切。海洋と大気を同量にせず、メタン生成を光合成にせず、紫外と赤外を取り違えず、土壌CO2による炭酸供給の風化促進を理解する。風化速度・炭素埋没の時間尺度も意識する。

\Needspace{4\baselineskip}\paragraph{選択肢の検討（査読待ち初稿）}\mbox{}\par

\noindent\textbf{1：不適切} 大気炭素全量は海洋と同程度。。海洋の炭素貯留は大気よりはるかに多く、同程度ではない。\par

\noindent\textbf{2：適切} 初期CO2は高く、減少に風化と光合成生物が関与。。地史的なCO2減少に風化と生物活動が関与する。光合成固定も呼吸で戻る部分と埋没する部分を区別する。\par

\noindent\textbf{3：不適切} 嫌気環境のメタン生成は光合成。。嫌気条件は整合するがメタン生成は光合成ではない。水素とCO2、酢酸等の基質を使う代謝である。\par

\noindent\textbf{4：不適切} CO2は紫外に幅広い吸収帯。。主要な温室作用は赤外領域の吸収・放射である。紫外との取り違え。\par

\noindent\textbf{5：不適切} 土壌CO2が高いので風化のCO2除去が減速。。土壌呼吸によるCO2供給は炭酸を作り化学的風化を促す。実際の速度は水流、表面積、温度、飽和度等にも依存する。\par

\paragraph{再利用する考え方} 物質の貯留量・輸送速度・化学形態を分け、系の境界と反応の前提を確かめる。貯留・反応・放射を区別する炭素循環を収支・平衡・反応速度のどのモデルで扱うか選ぶ。

\paragraph{共通理論} \hyperref[theory:geochemistry:carbon]{第\ref*{theory:geochemistry:carbon}節}, \hyperref[theory:geochemistry:weathering]{第\ref*{theory:geochemistry:weathering}節}, \hyperref[theory:geochemistry:redox]{第\ref*{theory:geochemistry:redox}節}

\par 参考：\cite{gc_carbon,atm_noaa_ghg,gc_redox}。

\paragraph{関連問題} \hyperref[q:h26-reg-applied-27]{H26-REG-APPLIED-27}

\par 難易度 3（校準初稿）、優先度未評価。

\section{二酸化炭素回収・貯留}

\subsection{平成26年度 問18}
\label{q:h26-reg-applied-18}

\noindent\textbf{原典：} sources/H26/17\_applied\_science.pdf、PDF物理7頁。

\textbf{公式正答：} 4（single）。

\paragraph{問題の要旨} CCSの回収と地中貯留に直接つながる技術・場所を比較する。

\Needspace{4\baselineskip}\paragraph{確認した条件・記述（独自要約）}\mbox{}\par

\noindent\textbf{方向：} 関係性が最も薄いもの。\par

\noindent\textbf{対象：} CCS（Carbon dioxide Capture and Storage）。\par

\paragraph{問われる核心} CO2回収・利用・地中貯留の工程

\paragraph{必要知識} CO2回収・利用・地中貯留の工程。CCSを回収と貯留に分けると、アミン・膜は回収、帯水層は貯留、CO2-EORは利用と貯留に接続する。界面活性剤4が相対的に最も関係が薄い。原典の「薄い」を「無関係」に強めず、EORの種類も区別する。

\paragraph{考え方・正答を選ぶ理由（査読待ち初稿）} CCSを回収と貯留に分けると、アミン・膜は回収、帯水層は貯留、CO2-EORは利用と貯留に接続する。界面活性剤4が相対的に最も関係が薄い。原典の「薄い」を「無関係」に強めず、EORの種類も区別する。

\Needspace{4\baselineskip}\paragraph{選択肢の検討（査読待ち初稿）}\mbox{}\par

\noindent\textbf{1：直接関係} 帯水層。深部塩水帯水層は貯留候補である。圧力、封じ込め、漏えい等を評価する。\par

\noindent\textbf{2：CO2-EORで関係} EOR。CO2を用いるEORは石油増進回収とCO2貯留に結び付く。他のEORまで同一ではない。\par

\noindent\textbf{3：直接関係} モノエタノールアミン吸収法。MEAの反応吸収はCO2回収の代表方式である。\par

\noindent\textbf{4：提示中で最も薄い} 界面活性剤。界面活性剤は油回収等に使われるが、提示中ではCO2分離・貯留そのものとの直接関係が薄い。全く関係がないと断定しない。\par

\noindent\textbf{5：直接関係} 膜分離。選択透過する膜でCO2を分離回収できる。\par

\paragraph{再利用する考え方} 物質の貯留量・輸送速度・化学形態を分け、系の境界と反応の前提を確かめる。CO2回収・利用・地中貯留の工程を収支・平衡・反応速度のどのモデルで扱うか選ぶ。

\paragraph{共通理論} \hyperref[theory:geochemistry:ccs]{第\ref*{theory:geochemistry:ccs}節}

\par 参考：\cite{gc_ccs,gc_capture,gc_mea}。

\par 難易度 3（校準初稿）、優先度未評価。

\section{地殻元素存在度と比較基準}

\subsection{平成25年度 問10}
\label{q:h25-reg-applied-10}

\noindent\textbf{原典：} sources/H25/17\_applied\_science.pdf、PDF物理5頁。

\textbf{公式正答：} 4（single）。

\paragraph{問題の要旨} 地殻元素の存在度の5組の大小を比較する。

\Needspace{4\baselineskip}\paragraph{確認した条件・記述（独自要約）}\mbox{}\par

\noindent\textbf{方向：} 最も不適切なもの。\par

\noindent\textbf{対象：} 地球表層の自然水・岩石・元素循環。各肢の条件は独自要約として保持。\par

\paragraph{問われる核心} 地殻組成・質量存在度と微量元素

\paragraph{必要知識} 地殻組成・質量存在度と微量元素。地殻という対象と質量存在度の基準を置くと、Cl>Kという4だけが逆である。大気・海水や地球全体の組成を代用しない。代表表の精密小数はモデル依存であるが、この大小比較は頑健である。

\paragraph{考え方・正答を選ぶ理由（査読待ち初稿）} 地殻という対象と質量存在度の基準を置くと、Cl>Kという4だけが逆である。大気・海水や地球全体の組成を代用しない。代表表の精密小数はモデル依存であるが、この大小比較は頑健である。

\Needspace{4\baselineskip}\paragraph{選択肢の検討（査読待ち初稿）}\mbox{}\par

\noindent\textbf{1：適切（通常条件）} O > Si。代表質量存在度46.60\%対27.72\%で成立。\par

\noindent\textbf{2：適切（通常条件）} Si > Ti。27.72\%対0.44\%で成立。\par

\noindent\textbf{3：適切（通常条件）} Al > Cu。Al8.13\%は微量のCuより多い。\par

\noindent\textbf{4：不適切} Cl > K。Cl0.020\%はK2.59\%よりはるかに少なく逆である。\par

\noindent\textbf{5：適切（通常条件）} Ti > Mn。Ti0.44\%はMn0.100\%より多い。\par

\paragraph{再利用する考え方} 物質の貯留量・輸送速度・化学形態を分け、系の境界と反応の前提を確かめる。地殻組成・質量存在度と微量元素を収支・平衡・反応速度のどのモデルで扱うか選ぶ。

\paragraph{共通理論} \hyperref[theory:geochemistry:crust]{第\ref*{theory:geochemistry:crust}節}

\par 参考：\cite{gc_crust,gc_hem}。

\par 難易度 3（校準初稿）、優先度未評価。

\section{溶質の移流・拡散分散と観測尺度}

\subsection{令和3年度 問18}
\label{q:r03-reg-applied-18}

\noindent\textbf{原典：} sources/R03/17\_applied\_science.pdf、PDF物理7頁。

\textbf{公式正答：} 5（single）。

\paragraph{問題の要旨} 分子拡散、機械分散、縦横の分散長と野外の尺度を比較する。

\Needspace{4\baselineskip}\paragraph{確認した条件・記述（独自要約）}\mbox{}\par

\noindent\textbf{方向：} 最も不適切なもの。\par

\noindent\textbf{対象：} 地球表層の自然水・岩石・元素循環。各肢の条件は独自要約として保持。\par

\paragraph{問われる核心} 移流拡散・機械分散と尺度依存

\paragraph{必要知識} 移流拡散・機械分散と尺度依存。5の野外がはるかに小さいという一般化が逆である。\(D_L=\alpha_L|v|+D_m^\ast\)と置けば、分子拡散係数は面積/時間、分散長は長さであり別の量である。流れの平均速度に比例する機械分散と静止状態にもある分子拡散を分け、現場尺度への外挿を検証する。

\paragraph{考え方・正答を選ぶ理由（査読待ち初稿）} 5の野外がはるかに小さいという一般化が逆である。\(D_L=\alpha_L|v|+D_m^\ast\)と置けば、分子拡散係数は面積/時間、分散長は長さであり別の量である。流れの平均速度に比例する機械分散と静止状態にもある分子拡散を分け、現場尺度への外挿を検証する。

\Needspace{4\baselineskip}\paragraph{選択肢の検討（査読待ち初稿）}\mbox{}\par

\noindent\textbf{1：適切（通常条件）} 分子の不規則運動が分子拡散。。不規則な熱運動に基づく濃度の平滑化である。\par

\noindent\textbf{2：適切（通常条件）} 濃度勾配に比例する速度で分子拡散。。Fick則は拡散流束\(J=-D_m\partial C/\partial x\)を与える。「速度」は流束として解釈し、正味の分子流速との同一視を避ける。\par

\noindent\textbf{3：適切（通常条件）} 不均質な間隙構造で流動中の溶質集団が広がる。。流速・経路の差を平均化した見かけの広がりで、静止水の分子拡散とは区別する。\par

\noindent\textbf{4：適切（通常条件）} 横分散長は通常縦より小さい。。流れ方向の速度差が大きく、通常は\(\alpha_T<\alpha_L\)である。媒体構造に依存する経験傾向。\par

\noindent\textbf{5：不適切} 野外分散長は室内よりはるかに小さい。。野外は大きな不均質性を含み一般に大きい側の傾向がある。1992年59現場の研究は尺度傾向と大きな散らばりを示し、普遍直線は否定する。\par

\paragraph{再利用する考え方} 物質の貯留量・輸送速度・化学形態を分け、系の境界と反応の前提を確かめる。移流拡散・機械分散と尺度依存を収支・平衡・反応速度のどのモデルで扱うか選ぶ。

\paragraph{共通理論} \hyperref[theory:geochemistry:transport]{第\ref*{theory:geochemistry:transport}節}

\par 参考：\cite{gc_dispersion,gc_hem}。

\par 難易度 3（校準初稿）、優先度未評価。

\section{天然水の電子受容体・酸化還元電位}

\subsection{令和4年度 問19}
\label{q:r04-reg-applied-19}

\noindent\textbf{原典：} sources/R04/17\_applied\_science.pdf、PDF物理6頁。

\textbf{公式正答：} 3（single）。

\par\textbf{検討未確定：} 肢4の酸素減少は酸化可能な物質と反応する条件、5は還元ガス・空気混合の条件が必要。原典の省略を無条件の保証として教えず部分校正。

\paragraph{問題の要旨} 電子授受、基準電極、有酸素水のEhと地下・火山水の状態を比較する。

\Needspace{4\baselineskip}\paragraph{確認した条件・記述（独自要約）}\mbox{}\par

\noindent\textbf{方向：} 最も不適切なもの。\par

\noindent\textbf{対象：} 地球表層の自然水・岩石・元素循環。各肢の条件は独自要約として保持。\par

\paragraph{問われる核心} 酸化還元電位・電極基準と自然水の電子受容体

\paragraph{必要知識} 酸化還元電位・電極基準と自然水の電子受容体。有酸素水のEhを-250mVとする3が不適切。Nernst式の酸素/水境界はpH7・空気下で約+805mVを与え、還元側の符号との違いを確かめられる。ただし実測Ehは複数反応・電極応答を含む。4と5の一般傾向は物質の供給・反応・混合条件を付して説明する。

\paragraph{考え方・正答を選ぶ理由（査読待ち初稿）} 有酸素水のEhを-250mVとする3が不適切。Nernst式の酸素/水境界はpH7・空気下で約+805mVを与え、還元側の符号との違いを確かめられる。ただし実測Ehは複数反応・電極応答を含む。4と5の一般傾向は物質の供給・反応・混合条件を付して説明する。

\Needspace{4\baselineskip}\paragraph{選択肢の検討（査読待ち初稿）}\mbox{}\par

\noindent\textbf{1：適切} 還元剤が電子供与、酸化剤が受容。。還元剤自身は酸化され、酸化剤自身は還元される。\par

\noindent\textbf{2：適切（基準指定）} 標準水素電極0を基準とした電位をEhと表す。。比較には同じ基準電極へ換算する。ORP計の表示を無換算でEhと混ぜない。\par

\noindent\textbf{3：不適切} pH7・25℃の有酸素水Ehは約-250mV。。酸素還元境界は\(E_h=1.229-0.05916\,\mathrm{pH}+0.01479\log_{10}p_{O_2}\)V（25℃、水活量1）。空気0.21・pH7で約+0.805V、-0.250Vは還元側であり方向が違う。実測値は混合電位で理想境界と一致しなくてもよい。\par

\noindent\textbf{4：反応条件を付す} 深部をゆっくり通る水はO2を失い還元状態へ。。長い滞留中の酸化可能な有機物や鉱物との反応でO2を消費し得る。単に時間が長いだけで反応物なしに酸素が消えるという保存則違反の説明にしない。\par

\noindent\textbf{5：通常の還元熱水・混合注意} 火山温泉湧水はSO2・H2S等と反応し普通はO2なし。。熱水は還元性ガス等と反応して酸素が乏しい場合が多い。湧出後の空気混合などを含め全温泉が永遠に無酸素とはしない。\par

\paragraph{再利用する考え方} 物質の貯留量・輸送速度・化学形態を分け、系の境界と反応の前提を確かめる。酸化還元電位・電極基準と自然水の電子受容体を収支・平衡・反応速度のどのモデルで扱うか選ぶ。

\paragraph{共通理論} \hyperref[theory:geochemistry:redox]{第\ref*{theory:geochemistry:redox}節}

\par 参考：\cite{gc_hem,gc_redox,gc_redox2009}。

\par 難易度 3（校準初稿）、優先度未評価。

\section{窒素固定・硝化脱窒と炭素源}

\subsection{令和2年度 問19}
\label{q:r02-reg-applied-19}

\noindent\textbf{原典：} sources/R02/17\_applied\_science.pdf、PDF物理7頁。

\textbf{公式正答：} 5（single）。

\par\textbf{検討未確定：} 公式5を保持。自養脱窒の一次実験例があるため肢5を全化学合成生物で否定できない。肢3も燃料窒素由来NOxとN2の新規固定を区別する範囲疑義。意図と科学的適用範囲を分けて査読待ち。

\paragraph{問題の要旨} 固定・硝化・脱窒と、無機炭素を使う独立栄養を区別する。

\Needspace{4\baselineskip}\paragraph{確認した条件・記述（独自要約）}\mbox{}\par

\noindent\textbf{方向：} 最も不適切なもの。\par

\noindent\textbf{対象：} 地球表層の自然水・岩石・元素循環。各肢の条件は独自要約として保持。\par

\paragraph{問われる核心} 窒素反応の電子収支と炭素源の独立な分類

\paragraph{必要知識} 窒素反応の電子収支と炭素源の独立な分類。公式5と初回5は一致するが、硝化と従属栄養脱窒という教科書的な対照だけでは5を普遍的誤りとできない。硫黄自養脱窒では無機電子供与体の酸化と硝酸還元がエネルギーを供給し、無機炭素を同化する。電子供与体・受容体・炭素源を別軸にして、原典と科学的例外を併記する。

\paragraph{考え方・正答を選ぶ理由（査読待ち初稿）} 公式5と初回5は一致するが、硝化と従属栄養脱窒という教科書的な対照だけでは5を普遍的誤りとできない。硫黄自養脱窒では無機電子供与体の酸化と硝酸還元がエネルギーを供給し、無機炭素を同化する。電子供与体・受容体・炭素源を別軸にして、原典と科学的例外を併記する。

\Needspace{4\baselineskip}\paragraph{選択肢の検討（査読待ち初稿）}\mbox{}\par

\noindent\textbf{1：適切（脱窒過程）} 嫌気硝酸呼吸で硝酸を還元しN2O・N2生成。。硝酸を電子受容体とする脱窒でN2O等を経てN2となる。中間生成物の量は条件で変わる。\par

\noindent\textbf{2：適切} アンモニア→亜硝酸→硝酸は硝化。。酸化数を上げる硝化の二段階である。代表の独立栄養硝化生物は無機酸化のエネルギーで炭素を固定する。\par

\noindent\textbf{3：固定源を分ける条件付き} 雷・燃焼NOxは非生物的窒素固定。。雷や高温の熱的NOxはN2から反応性窒素を作る非生物固定である。ただし燃料に既に含まれた窒素の放出まで全て新規N2固定とすることはできない。\par

\noindent\textbf{4：適切（還元と同化）} 固定細菌・藍藻はN2からアミノ酸合成。。窒素固定ではN2をまず還元してアンモニアを得て、アミノ酸へ同化する。生成物と最初の酵素反応を区別する。\par

\noindent\textbf{5：意図上不適切・自養脱窒の例外} 化学合成生物は脱窒のエネルギーでCO2から有機物合成。。従属栄養脱窒を念頭にした設問意図では硝化との混同である。しかし1989年一次研究は嫌気条件で硫黄を電子供与体、HCO3-を炭素源とする自養脱窒菌を集積しているので、脱窒と無機炭素固定の結合は実在し一律に否定できない。\par

\paragraph{再利用する考え方} 物質の貯留量・輸送速度・化学形態を分け、系の境界と反応の前提を確かめる。窒素反応の電子収支と炭素源の独立な分類を収支・平衡・反応速度のどのモデルで扱うか選ぶ。

\paragraph{共通理論} \hyperref[theory:geochemistry:nitrogen]{第\ref*{theory:geochemistry:nitrogen}節}, \hyperref[theory:geochemistry:redox]{第\ref*{theory:geochemistry:redox}節}

\par 参考：\cite{gc_sulfur_denit,gc_redox,gc_hem}。

\par 難易度 3（校準初稿）、優先度未評価。

\section{陸水の主要イオン・地質と採水変化}

\subsection{平成24年度 問33}
\label{q:h24-reg-applied-33}

\noindent\textbf{原典：} sources/H24/17\_applied\_science.pdf、PDF物理10頁。

\textbf{公式正答：} 5（single）。

\paragraph{問題の要旨} 地下水の鉄沈殿、鉱山酸性水、日本のケイ酸、採水pHと世界河川のイオン源を判定する。

\Needspace{4\baselineskip}\paragraph{確認した条件・記述（独自要約）}\mbox{}\par

\noindent\textbf{方向：} 最も不適切なもの。\par

\noindent\textbf{対象：} 地球表層の自然水・岩石・元素循環。各肢の条件は独自要約として保持。\par

\paragraph{問われる核心} 陸水の組成を決める供給源と採水後の反応

\paragraph{必要知識} 陸水の組成を決める供給源と採水後の反応。風化、酸化還元、脱気、海塩を別の過程として読む。岩石との反応によるCa-HCO3系を落としてNaCl起源に限定する5が不適切である。日本と世界のケイ酸濃度は集計時期・SiとSiO2の表示基準を添える。

\paragraph{考え方・正答を選ぶ理由（査読待ち初稿）} 風化、酸化還元、脱気、海塩を別の過程として読む。岩石との反応によるCa-HCO3系を落としてNaCl起源に限定する5が不適切である。日本と世界のケイ酸濃度は集計時期・SiとSiO2の表示基準を添える。

\Needspace{4\baselineskip}\paragraph{選択肢の検討（査読待ち初稿）}\mbox{}\par

\noindent\textbf{1：適切（通常条件）} 低酸素地下水の湧出に伴い鉄が酸化して茶褐色沈殿。。還元的な地下でFe(II)が溶存し、湧出後に酸素を得るとFe(III)の水酸化物等が沈殿する。茶褐色沈殿を地下の酸化状態と同一視しない。\par

\noindent\textbf{2：適切（通常条件）} 硫化鉱物酸化で硫酸増加・pH低下。。硫化鉱物の酸化は硫酸と酸を生成する。炭酸塩等の中和能でpHの結果は変わるが、酸性鉱山排水の典型的原因である。\par

\noindent\textbf{3：適切（通常条件）} 日本の平均河川は世界より溶存ケイ酸が高い。。1970年一次論文の表4は、日本の蒸発残渣74.8mg/L・SiO2比26.84\%、世界100mg/L・11.67\%であり、濃度へ戻すと約20.1対11.7mg/Lである。古い集計の傾向であって全河川の普遍命題ではない。\par

\noindent\textbf{4：適切（通常条件）} 現場の非攪拌水は室内放置水よりpH低い傾向。。地下水は土壌CO2を多く溶かしている場合があり、採取・攪拌による脱気で炭酸平衡が動きpHが上がる。CO2過飽和という条件を付し、鉄酸化等で逆の変化もあり得るため現場測定を優先する。\par

\noindent\textbf{5：不適切} 世界河川の大半はNaCl主成分・海塩が主供給源。。河川の平均的な主要イオンはCa・HCO3等の岩石風化に依存する。海塩NaClは雨・沿岸・混合で重要だが世界河川の大半を一律に海塩支配とできない。\par

\paragraph{再利用する考え方} 物質の貯留量・輸送速度・化学形態を分け、系の境界と反応の前提を確かめる。陸水の組成を決める供給源と採水後の反応を収支・平衡・反応速度のどのモデルで扱うか選ぶ。

\paragraph{共通理論} \hyperref[theory:geochemistry:river]{第\ref*{theory:geochemistry:river}節}, \hyperref[theory:geochemistry:redox]{第\ref*{theory:geochemistry:redox}節}, \hyperref[theory:geochemistry:carbonate]{第\ref*{theory:geochemistry:carbonate}節}, \hyperref[theory:geochemistry:sampling]{第\ref*{theory:geochemistry:sampling}節}

\par 参考：\cite{gc_river1970,gc_hem,gc_redox,gc_mine}。

\paragraph{関連問題} \hyperref[q:h29-reg-applied-20]{H29-REG-APPLIED-20}

\par 難易度 3（校準初稿）、優先度未評価。

\subsection{平成25年度 問31}
\label{q:h25-reg-applied-31}

\noindent\textbf{原典：} sources/H25/17\_applied\_science.pdf、PDF物理11頁。

\textbf{公式正答：} 4（single）。

\par\textbf{検討未確定：} 原典は「炭酸カルシウム型」である。通常の主要溶存イオンによる名称はCa-HCO3（炭酸水素カルシウム）型で、CaCO3分子が主成分との理解は誤る。原典の名称を保持して部分校正。

\paragraph{問題の要旨} 希薄化、石灰岩反応、蒸発濃縮、水素酸素同位体と酸性河川の起源を比較する。

\Needspace{4\baselineskip}\paragraph{確認した条件・記述（独自要約）}\mbox{}\par

\noindent\textbf{方向：} 最も不適切なもの。\par

\noindent\textbf{対象：} 地球表層の自然水・岩石・元素循環。各肢の条件は独自要約として保持。\par

\paragraph{問われる核心} 陸水の反応・濃縮と同位体による水の起源

\paragraph{必要知識} 陸水の反応・濃縮と同位体による水の起源。重い同位体の分別がHとOで連動するため、天水線は正の相関を持つ。反比例とする4が不適切である。Ca-HCO3型の名称、蒸発濃縮と海水混合の違いを補足し、原典の全記述を無条件の一般法則へしない。

\paragraph{考え方・正答を選ぶ理由（査読待ち初稿）} 重い同位体の分別がHとOで連動するため、天水線は正の相関を持つ。反比例とする4が不適切である。Ca-HCO3型の名称、蒸発濃縮と海水混合の違いを補足し、原典の全記述を無条件の一般法則へしない。

\Needspace{4\baselineskip}\paragraph{選択肢の検討（査読待ち初稿）}\mbox{}\par

\noindent\textbf{1：適切（一般傾向）} 日本は多雨急流・低蒸発濃縮で成分が一般に希薄。。古い平均比較は支持する。流域地質・汚染・温泉混合等により例外がある。\par

\noindent\textbf{2：溶解は適切・名称疑義} 石灰岩溶解で弱アルカリ性・Ca増加・炭酸カルシウム型。。炭酸を含む水に石灰岩が溶けるとCaとHCO3が増す。「炭酸カルシウム型」という原本の呼び方と溶存種を区別する。\par

\noindent\textbf{3：適切（濃縮混合傾向）} 乾燥蒸発や河口で塩が濃くなり海水組成へ近づく傾向。。濃縮・海水混合は塩類濃度を上げる。ただし蒸発だけで全イオン比が海水へ必ず一致するわけではない。\par

\noindent\textbf{4：不適切} HとOの同位体比はおおむね反比例。。天水のδ2Hとδ18Oは正の相関で、代表線はδ2H=\allowbreak{}8δ18O+10。反比例ではない。海水起源の水蒸気が分別するという説明と区別する。\par

\noindent\textbf{5：適切} 火山・酸性降雨・鉱山排水が酸性水源。。これらは酸を供給する代表過程である。pHの結果は流域の中和反応にも依存する。\par

\paragraph{再利用する考え方} 物質の貯留量・輸送速度・化学形態を分け、系の境界と反応の前提を確かめる。陸水の反応・濃縮と同位体による水の起源を収支・平衡・反応速度のどのモデルで扱うか選ぶ。

\paragraph{共通理論} \hyperref[theory:geochemistry:river]{第\ref*{theory:geochemistry:river}節}, \hyperref[theory:geochemistry:carbonate]{第\ref*{theory:geochemistry:carbonate}節}, \hyperref[theory:geochemistry:isotope]{第\ref*{theory:geochemistry:isotope}節}

\par 参考：\cite{gc_river1970,gc_hem,gc_meteoric}。

\paragraph{関連問題} \hyperref[q:r01-reg-applied-20]{R01-REG-APPLIED-20}

\par 難易度 3（校準初稿）、優先度未評価。

\subsection{平成27年度 問27}
\label{q:h27-reg-applied-27}

\noindent\textbf{原典：} sources/H27/17\_applied\_science.pdf、PDF物理10頁。

\textbf{公式正答：} 5（single）。

\par\textbf{独立解答と公式正答に相違あり：} 独立解答 4|5。初回4|5と公式5を保持。肢2はCaと酸化条件、4は海洋逆風化条件、5は成長による取り込みと死後殻溶解の区別を省略。「ことがある」を原典通り保持し一意な科学的誤り肢の確定は査読待ち。

\paragraph{問題の要旨} ケイ酸供給、生物取り込み、硫化物反応、還元地下水と海洋反応を比較する。

\Needspace{4\baselineskip}\paragraph{確認した条件・記述（独自要約）}\mbox{}\par

\noindent\textbf{方向：} 最も不適切なもの。\par

\noindent\textbf{対象：} 地球表層の自然水・岩石・元素循環。各肢の条件は独自要約として保持。\par

\paragraph{問われる核心} 生物によるケイ酸取り込みと溶解・逆風化

\paragraph{必要知識} 生物によるケイ酸取り込みと溶解・逆風化。公式は5だが初回は4|5であった。珪藻成長では水からケイ酸を除くため意図上の誤答理由は説明できる。一方、死後殻の溶解放出と海洋逆風化の反応条件まで含めると、原典の省略された条件が重要である。公式から原典5を「成長時に」と改変せず、各肢の条件と未解決を保つ。

\paragraph{考え方・正答を選ぶ理由（査読待ち初稿）} 公式は5だが初回は4|5であった。珪藻成長では水からケイ酸を除くため意図上の誤答理由は説明できる。一方、死後殻の溶解放出と海洋逆風化の反応条件まで含めると、原典の省略された条件が重要である。公式から原典5を「成長時に」と改変せず、各肢の条件と未解決を保つ。

\Needspace{4\baselineskip}\paragraph{選択肢の検討（査読待ち初稿）}\mbox{}\par

\noindent\textbf{1：適切（集計傾向）} 日本の平均河川は世界より溶存ケイ酸が高い。。1970年一次論文の表4は、日本の蒸発残渣74.8mg/L・SiO2比26.84\%、世界100mg/L・11.67\%であり、濃度へ戻すと約20.1対11.7mg/Lである。古い集計の傾向であって全河川の普遍命題ではない。\par

\noindent\textbf{2：条件付き・省略疑義} 硫化物に富む底質なら硫酸カルシウム型河川水。。硫化物が酸化されれば硫酸が生じ、酸によりCa鉱物が溶ければCa-SO4型を形成し得る。ただし硫化物の存在だけでは酸素・Ca供給を保証しない。\par

\noindent\textbf{3：適切（反応条件）} 低酸素地下水は酸化反応でO2を消耗し還元的。。供給O2が少なく還元性成分と反応すると地下水の還元状態が進み、FeMn等の形態が変わる。\par

\noindent\textbf{4：逆風化条件で成立・省略疑義} アルミノケイ酸塩の海洋供給で海水pH低下。。海洋の粘土生成（逆風化）はアルカリ度を消費してCO2を生成し、pHを下げ得る。陸上のケイ酸塩風化による酸消費とは別過程である。原典が反応条件を省いており供給だけで必ず低下とは断定できない。\par

\noindent\textbf{5：意図上不適切・生活史範囲疑義} ケイ藻でケイ素が放出され河川濃度が上がる場合。。生きた珪藻はケイ酸を取り込み殻を作るので、成長を濃度上昇の機構とするなら逆である。ただし死後殻の溶解はケイ酸を再供給し得るため、「ことがある」という原典の広い表現は一律に誤りと確定できない。\par

\paragraph{再利用する考え方} 物質の貯留量・輸送速度・化学形態を分け、系の境界と反応の前提を確かめる。生物によるケイ酸取り込みと溶解・逆風化を収支・平衡・反応速度のどのモデルで扱うか選ぶ。

\paragraph{共通理論} \hyperref[theory:geochemistry:silica]{第\ref*{theory:geochemistry:silica}節}, \hyperref[theory:geochemistry:river]{第\ref*{theory:geochemistry:river}節}, \hyperref[theory:geochemistry:redox]{第\ref*{theory:geochemistry:redox}節}

\par 参考：\cite{gc_river1970,gc_hem,gc_diatom,gc_reverse}。

\paragraph{関連問題} \hyperref[q:r01-reg-applied-20]{R01-REG-APPLIED-20}

\par 難易度 3（校準初稿）、優先度未評価。

\subsection{平成29年度 問20}
\label{q:h29-reg-applied-20}

\noindent\textbf{原典：} sources/H29/17\_applied\_science.pdf、PDF物理6頁。

\textbf{公式正答：} 1（single）。

\paragraph{問題の要旨} 海塩支配の一般化、硫化物酸化、ケイ酸濃度、脱気pH、鉄沈殿を比較する。

\Needspace{4\baselineskip}\paragraph{確認した条件・記述（独自要約）}\mbox{}\par

\noindent\textbf{方向：} 最も不適切なもの。\par

\noindent\textbf{対象：} 地球表層の自然水・岩石・元素循環。各肢の条件は独自要約として保持。\par

\paragraph{問われる核心} 主要イオンの供給と採水後の化学反応

\paragraph{必要知識} 主要イオンの供給と採水後の化学反応。風化由来Ca-HCO3等を無視して世界の河川を海塩NaClに限定する1が不適切である。採水後のpH変化を地質の違いと混ぜず、CO2脱気の条件として考える。鉄は地下の溶存Fe(II)と地表沈殿Fe(III)を区別する。

\paragraph{考え方・正答を選ぶ理由（査読待ち初稿）} 風化由来Ca-HCO3等を無視して世界の河川を海塩NaClに限定する1が不適切である。採水後のpH変化を地質の違いと混ぜず、CO2脱気の条件として考える。鉄は地下の溶存Fe(II)と地表沈殿Fe(III)を区別する。

\Needspace{4\baselineskip}\paragraph{選択肢の検討（査読待ち初稿）}\mbox{}\par

\noindent\textbf{1：不適切} 世界河川の大半はNaCl主成分・海塩起源。。河川の平均的な主要イオンはCa・HCO3等の岩石風化に依存する。海塩NaClは雨・沿岸・混合で重要だが世界河川の大半を一律に海塩支配とできない。\par

\noindent\textbf{2：適切（通常条件）} 硫化鉱物酸化で硫酸増・pH低下。。硫化鉱物の酸化は硫酸と酸を生成する。炭酸塩等の中和能でpHの結果は変わるが、酸性鉱山排水の典型的原因である。\par

\noindent\textbf{3：適切（通常条件）} 日本平均河川ケイ酸は世界より高い。。1970年一次論文の表4は、日本の蒸発残渣74.8mg/L・SiO2比26.84\%、世界100mg/L・11.67\%であり、濃度へ戻すと約20.1対11.7mg/Lである。古い集計の傾向であって全河川の普遍命題ではない。\par

\noindent\textbf{4：適切（通常条件）} 非攪拌地下水pHは攪拌水より低い傾向。。地下水は土壌CO2を多く溶かしている場合があり、採取・攪拌による脱気で炭酸平衡が動きpHが上がる。CO2過飽和という条件を付し、鉄酸化等で逆の変化もあり得るため現場測定を優先する。\par

\noindent\textbf{5：適切（通常条件）} 低酸素地下水湧出で鉄の酸化沈殿。。還元的な地下でFe(II)が溶存し、湧出後に酸素を得るとFe(III)の水酸化物等が沈殿する。茶褐色沈殿を地下の酸化状態と同一視しない。\par

\paragraph{再利用する考え方} 物質の貯留量・輸送速度・化学形態を分け、系の境界と反応の前提を確かめる。主要イオンの供給と採水後の化学反応を収支・平衡・反応速度のどのモデルで扱うか選ぶ。

\paragraph{共通理論} \hyperref[theory:geochemistry:river]{第\ref*{theory:geochemistry:river}節}, \hyperref[theory:geochemistry:carbonate]{第\ref*{theory:geochemistry:carbonate}節}, \hyperref[theory:geochemistry:sampling]{第\ref*{theory:geochemistry:sampling}節}, \hyperref[theory:geochemistry:redox]{第\ref*{theory:geochemistry:redox}節}

\par 参考：\cite{gc_river1970,gc_hem,gc_redox,gc_mine}。

\paragraph{関連問題} \hyperref[q:h24-reg-applied-33]{H24-REG-APPLIED-33}

\par 難易度 3（校準初稿）、優先度未評価。

\subsection{令和元年度 問20}
\label{q:r01-reg-applied-20}

\noindent\textbf{原典：} sources/R01/17\_applied\_science.pdf、PDF物理6頁。

\textbf{公式正答：} 1（single）。

\par\textbf{検討未確定：} 原典は「炭酸カルシウム型」である。通常の主要溶存イオンによる名称はCa-HCO3（炭酸水素カルシウム）型で、CaCO3分子が主成分との理解は誤る。原典の名称を保持して部分校正。

\paragraph{問題の要旨} FeMnの溶存、火山地質ケイ酸、ORP、石灰岩反応と同位体緯度差を比較する。

\Needspace{4\baselineskip}\paragraph{確認した条件・記述（独自要約）}\mbox{}\par

\noindent\textbf{方向：} 最も不適切なもの。\par

\noindent\textbf{対象：} 地球表層の自然水・岩石・元素循環。各肢の条件は独自要約として保持。\par

\paragraph{問われる核心} 自然水の酸化還元・地質と溶存組成

\paragraph{必要知識} 自然水の酸化還元・地質と溶存組成。溶存FeMnを考えると河川と還元地下水の大小を逆にした1が不適切である。粒子を含む全濃度と溶存濃度も区別する。4の水質名称は原典を保持して疑義とし、炭酸平衡の溶存種はCa2+・HCO3-と説明する。

\paragraph{考え方・正答を選ぶ理由（査読待ち初稿）} 溶存FeMnを考えると河川と還元地下水の大小を逆にした1が不適切である。粒子を含む全濃度と溶存濃度も区別する。4の水質名称は原典を保持して疑義とし、炭酸平衡の溶存種はCa2+・HCO3-と説明する。

\Needspace{4\baselineskip}\paragraph{選択肢の検討（査読待ち初稿）}\mbox{}\par

\noindent\textbf{1：不適切} 河川水は地下水よりFeMnに富むことが多い。。酸化的河川ではFeMn酸化物が析出し、還元地下水はFe2+・Mn2+が溶存し得る。一般の大小方向が逆である。\par

\noindent\textbf{2：適切（地質傾向）} 新生代火山系を流れる河川でケイ酸が高い傾向。。火山性の反応しやすい物質からのケイ酸供給は高濃度へ寄与し得る。地質だけで全流域の濃度を固定しない。\par

\noindent\textbf{3：適切（基準と条件）} 大気共存水のORPは一般に正。。酸素が供給される通常の天然水は酸化的で、電位は正側の傾向。電極基準・温度pH・混合電位を区別する。\par

\noindent\textbf{4：反応は適切・名称疑義} 石灰岩溶解でCa増・弱アルカリ・炭酸カルシウム型。。石灰岩と炭酸の反応でCa・HCO3が増加する。原典の炭酸カルシウム型と通常のCa-HCO3名称を区別する。\par

\noindent\textbf{5：適切（一般傾向）} 降水酸素同位体値は赤道側で大きく極へ小さい。。水蒸気の逐次凝結で重い同位体が除かれ、高緯度降水は相対的に低いδ18Oとなる一般傾向。局地気候等の例外もある。\par

\paragraph{再利用する考え方} 物質の貯留量・輸送速度・化学形態を分け、系の境界と反応の前提を確かめる。自然水の酸化還元・地質と溶存組成を収支・平衡・反応速度のどのモデルで扱うか選ぶ。

\paragraph{共通理論} \hyperref[theory:geochemistry:river]{第\ref*{theory:geochemistry:river}節}, \hyperref[theory:geochemistry:redox]{第\ref*{theory:geochemistry:redox}節}, \hyperref[theory:geochemistry:carbonate]{第\ref*{theory:geochemistry:carbonate}節}, \hyperref[theory:geochemistry:isotope]{第\ref*{theory:geochemistry:isotope}節}

\par 参考：\cite{gc_hem,gc_river1970,gc_redox,gc_meteoric}。

\paragraph{関連問題} \hyperref[q:h25-reg-applied-31]{H25-REG-APPLIED-31}, \hyperref[q:h27-reg-applied-27]{H27-REG-APPLIED-27}

\par 難易度 3（校準初稿）、優先度未評価。

\section{水収支・貯留・平均滞留時間}

\subsection{平成23年度 問22}
\label{q:h23-reg-applied-22}

\noindent\textbf{原典：} sources/H23/17\_applied\_science.pdf、PDF物理7頁。

\textbf{公式正答：} 4（single）。

\paragraph{問題の要旨} 提示された陸上の水貯留の中で最大のものを選ぶ。

\Needspace{4\baselineskip}\paragraph{確認した条件・記述（独自要約）}\mbox{}\par

\noindent\textbf{方向：} 最も貯留量が大きいもの。\par

\noindent\textbf{対象：} 湖沼・河川・湿地・地下・土壌の5貯留。\par

\paragraph{問われる核心} 水の貯留量と対象集合

\paragraph{必要知識} 水の貯留量と対象集合。示された5形態では地下水4が最大である。氷河や海洋は候補にない。水の量をストックとして比較し、流量や取り出せる資源量のランキングへ移さない。

\paragraph{考え方・正答を選ぶ理由（査読待ち初稿）} 示された5形態では地下水4が最大である。氷河や海洋は候補にない。水の量をストックとして比較し、流量や取り出せる資源量のランキングへ移さない。

\Needspace{4\baselineskip}\paragraph{選択肢の検討（査読待ち初稿）}\mbox{}\par

\noindent\textbf{1：最大でない} 湖沼水。USGS表の全湖沼は約17.6万km3で、地下水より小さい。\par

\noindent\textbf{2：最大でない} 河川水。河川の瞬時貯留は約2120km3である。大きい年間輸送量と区別する。\par

\noindent\textbf{3：最大でない} 湿地の水。湿地の貯留は約11470km3で地下水より小さい。\par

\noindent\textbf{4：最大} 地下水。全地下水約2340万km3、淡水地下水だけでも約1053万km3で提示中最大。\par

\noindent\textbf{5：最大でない} 土壌水。土壌水約16500km3は地下水より小さい。\par

\paragraph{再利用する考え方} 物質の貯留量・輸送速度・化学形態を分け、系の境界と反応の前提を確かめる。水の貯留量と対象集合を収支・平衡・反応速度のどのモデルで扱うか選ぶ。

\paragraph{共通理論} \hyperref[theory:geochemistry:reservoir]{第\ref*{theory:geochemistry:reservoir}節}

\par 参考：\cite{gc_water}。

\paragraph{関連問題} \hyperref[q:h26-reg-applied-25]{H26-REG-APPLIED-25}

\par 難易度 3（校準初稿）、優先度未評価。

\subsection{平成26年度 問25}
\label{q:h26-reg-applied-25}

\noindent\textbf{原典：} sources/H26/17\_applied\_science.pdf、PDF物理11頁。

\textbf{公式正答：} 5（single）。

\par\textbf{検討未確定：} 原典の陸水に対する雪氷約70\%は、淡水貯留を分母にすれば68.7\%という代表表と合う。海洋以外の全水（塩性地下水等を含む）を陸水とすると約50\%に変わるため分母疑義を保持。

\paragraph{問題の要旨} 全球水量の概数・分母と、河川水と大気水の貯留量を比較する。

\Needspace{4\baselineskip}\paragraph{確認した条件・記述（独自要約）}\mbox{}\par

\noindent\textbf{方向：} 最も不適切なもの。\par

\noindent\textbf{対象：} 地球表層の自然水・岩石・元素循環。各肢の条件は独自要約として保持。\par

\paragraph{問われる核心} 全球水貯留・概数と割合の分母

\paragraph{必要知識} 全球水貯留・概数と割合の分母。同じ貯留量で比較すると河川2120と大気12900km3であり、多いとする5が不適切である。他の割合は丸めと分母を確認する。雪氷約70\%は淡水分母の説明として示し、原典の陸水という広い語は疑義を残す。

\paragraph{考え方・正答を選ぶ理由（査読待ち初稿）} 同じ貯留量で比較すると河川2120と大気12900km3であり、多いとする5が不適切である。他の割合は丸めと分母を確認する。雪氷約70\%は淡水分母の説明として示し、原典の陸水という広い語は疑義を残す。

\Needspace{4\baselineskip}\paragraph{選択肢の検討（査読待ち初稿）}\mbox{}\par

\noindent\textbf{1：適切（概数）} 総量約14億km3。1386000000km3という代表表を約14億km3へ丸めており妥当。\par

\noindent\textbf{2：適切（概数）} 海水約97\%。海水96.54\%を約97\%へ丸めている。\par

\noindent\textbf{3：分母疑義} 陸水の雪氷約70\%。淡水分母の氷は68.7\%で約70\%。陸水に塩性地下水まで含めるなら同じ割合でない。\par

\noindent\textbf{4：適切} 湖沼は地下水より少ない。。湖沼全体176400km3は全地下水23400000km3より小さい。\par

\noindent\textbf{5：不適切} 河川水は大気水蒸気より多い。。河川2120km3は大気12900km3より小さい。年間流出量をストックへ代入しない。\par

\paragraph{再利用する考え方} 物質の貯留量・輸送速度・化学形態を分け、系の境界と反応の前提を確かめる。全球水貯留・概数と割合の分母を収支・平衡・反応速度のどのモデルで扱うか選ぶ。

\paragraph{共通理論} \hyperref[theory:geochemistry:reservoir]{第\ref*{theory:geochemistry:reservoir}節}

\par 参考：\cite{gc_water}。

\paragraph{関連問題} \hyperref[q:h23-reg-applied-22]{H23-REG-APPLIED-22}, \hyperref[q:r01-ret-applied-18]{R01-RET-APPLIED-18}

\par 難易度 3（校準初稿）、優先度未評価。

\subsection{令和元年度再試験 問18}
\label{q:r01-ret-applied-18}

\noindent\textbf{原典：} sources/R01\_retest/17\_applied\_science.pdf、PDF物理6頁。

\textbf{公式正答：} 4（single）。

\par\textbf{検討未確定：} 原典の陸水に対する雪氷約70\%は、淡水貯留を分母にすれば68.7\%という代表表と合う。海洋以外の全水（塩性地下水等を含む）を陸水とすると約50\%に変わるため分母疑義を保持。

\paragraph{問題の要旨} 全球水の割合と地下・湖・河・大気の貯留を比較する。

\Needspace{4\baselineskip}\paragraph{確認した条件・記述（独自要約）}\mbox{}\par

\noindent\textbf{方向：} 最も不適切なもの。\par

\noindent\textbf{対象：} 地球表層の自然水・岩石・元素循環。各肢の条件は独自要約として保持。\par

\paragraph{問われる核心} 水貯留の分母と湖沼の淡水限定

\paragraph{必要知識} 水貯留の分母と湖沼の淡水限定。河川2120km3は大気12900km3より小さいため4が不適切。H26問25から候補順が変わり、湖沼に淡水という限定もあるので番号を転用しない。雪氷70\%の分母は淡水なら整合するが陸水という原語の範囲疑義を残す。

\paragraph{考え方・正答を選ぶ理由（査読待ち初稿）} 河川2120km3は大気12900km3より小さいため4が不適切。H26問25から候補順が変わり、湖沼に淡水という限定もあるので番号を転用しない。雪氷70\%の分母は淡水なら整合するが陸水という原語の範囲疑義を残す。

\Needspace{4\baselineskip}\paragraph{選択肢の検討（査読待ち初稿）}\mbox{}\par

\noindent\textbf{1：適切（概数）} 総量約14億km3。1386000000km3という代表表を約14億km3へ丸めており妥当。\par

\noindent\textbf{2：適切（概数）} 海水約97\%。海水96.54\%を約97\%へ丸めている。\par

\noindent\textbf{3：分母疑義} 陸水中雪氷約70\%。淡水分母の氷は68.7\%で約70\%。陸水に塩性地下水まで含めるなら同じ割合でない。\par

\noindent\textbf{4：不適切} 河川水は大気の水蒸気より多い。。河川2120km3は大気12900km3より小さい。年間流出量をストックへ代入しない。\par

\noindent\textbf{5：適切} 淡水湖は地下水より少ない。。淡水湖91000km3は淡水地下水10530000km3より小さく、全地下水と比較しても小さい。\par

\paragraph{再利用する考え方} 物質の貯留量・輸送速度・化学形態を分け、系の境界と反応の前提を確かめる。水貯留の分母と湖沼の淡水限定を収支・平衡・反応速度のどのモデルで扱うか選ぶ。

\paragraph{共通理論} \hyperref[theory:geochemistry:reservoir]{第\ref*{theory:geochemistry:reservoir}節}

\par 参考：\cite{gc_water}。

\paragraph{関連問題} \hyperref[q:h26-reg-applied-25]{H26-REG-APPLIED-25}

\par 難易度 3（校準初稿）、優先度未評価。

\section{天然水の電気伝導率・温度と組成}

\subsection{令和3年度 問19}
\label{q:r03-reg-applied-19}

\noindent\textbf{原典：} sources/R03/17\_applied\_science.pdf、PDF物理7頁。

\textbf{公式正答：} 2（single）。

\paragraph{問題の要旨} 天然水のイオン量・温度と電気伝導率、抵抗率と残留物を比較する。

\Needspace{4\baselineskip}\paragraph{確認した条件・記述（独自要約）}\mbox{}\par

\noindent\textbf{方向：} 最も不適切なもの。\par

\noindent\textbf{対象：} 地球表層の自然水・岩石・元素循環。各肢の条件は独自要約として保持。\par

\paragraph{問われる核心} イオンの伝導・抵抗率と温度補正

\paragraph{必要知識} イオンの伝導・抵抗率と温度補正。通常の天然水の電気伝導率は温度上昇で増すので2が不適切。\(\kappa=\sum_i \lambda_i c_i\)という希薄な加算近似では、イオン種別の伝導寄与と濃度の双方が必要である。比較時は温度を揃え、25℃換算値か実測温度の値かを明記する。

\paragraph{考え方・正答を選ぶ理由（査読待ち初稿）} 通常の天然水の電気伝導率は温度上昇で増すので2が不適切。\(\kappa=\sum_i \lambda_i c_i\)という希薄な加算近似では、イオン種別の伝導寄与と濃度の双方が必要である。比較時は温度を揃え、25℃換算値か実測温度の値かを明記する。

\Needspace{4\baselineskip}\paragraph{選択肢の検討（査読待ち初稿）}\mbox{}\par

\noindent\textbf{1：適切（通常条件）} 伝導率で溶存イオン量を大まかに比較。。イオン濃度と移動度に依存するので概略指標になる。中性有機物等の全溶存量の直接測定ではない。\par

\noindent\textbf{2：不適切} 温度が上がると伝導率低下。。同一試料・通常温度では温度上昇でイオン移動度が増し伝導率は上がる。25℃への温度補正値と生の温度依存を分ける。\par

\noindent\textbf{3：適切（通常条件）} 通常海水→河川→雨の順で低下。。塩類濃度の一般傾向に沿う。汚染・温泉・酸性水等の例外があり試料名だけで固定しない。\par

\noindent\textbf{4：適切（通常条件）} 1cm立方体の対面間に対応する電気抵抗率の逆数。。伝導率\(\kappa=1/\rho=\ell/(RA)\)。\(\ell=1\,\mathrm{cm},A=1\,\mathrm{cm}^2\)で測定抵抗Rの逆数の数値に対応するが単位はS/cm。抵抗と抵抗率の量を混ぜない。\par

\noindent\textbf{5：適切（通常条件）} 同じ組成比で濃度が違えば蒸発残留物と密接に関係。。同じ組成・温度・希薄条件なら比例的に変化する。ただし異なるイオン組成の試料を一つの係数で換算しない。\par

\paragraph{再利用する考え方} 物質の貯留量・輸送速度・化学形態を分け、系の境界と反応の前提を確かめる。イオンの伝導・抵抗率と温度補正を収支・平衡・反応速度のどのモデルで扱うか選ぶ。

\paragraph{共通理論} \hyperref[theory:geochemistry:conductivity]{第\ref*{theory:geochemistry:conductivity}節}

\par 参考：\cite{gc_hem,gc_conductance,gc_conductance1994}。

\par 難易度 3（校準初稿）、優先度未評価。

\section{鉱物水反応・活量と風化生成物}

\subsection{令和4年度 問20}
\label{q:r04-reg-applied-20}

\noindent\textbf{原典：} sources/R04/17\_applied\_science.pdf、PDF物理7頁。

\textbf{公式正答：} 1（single）。

\par\textbf{検討未確定：} 肢3は残存固相・固定条件での飽和平衡を必要とする。少量全溶解・開放流動を含めれば原典の無条件な最終平衡説明は成立しないため条件疑義を保持。

\paragraph{問題の要旨} 活量係数の理論、イオン強度、溶解平衡、風化系列と速度を区別する。

\Needspace{4\baselineskip}\paragraph{確認した条件・記述（独自要約）}\mbox{}\par

\noindent\textbf{方向：} 最も不適切なもの。\par

\noindent\textbf{対象：} 地球表層の自然水・岩石・元素循環。各肢の条件は独自要約として保持。\par

\paragraph{問われる核心} 溶解の平衡・活量と反応速度の分離

\paragraph{必要知識} 溶解の平衡・活量と反応速度の分離。1は理論の対象を取り違えている。Debye-Huckelはイオン強度から活量係数を近似する理論で、活量による平衡計算へ接続する。平衡でどこまで溶けるかと、そこへ至る速度、風化抵抗という経験系列を別の概念として整理する。

\paragraph{考え方・正答を選ぶ理由（査読待ち初稿）} 1は理論の対象を取り違えている。Debye-Huckelはイオン強度から活量係数を近似する理論で、活量による平衡計算へ接続する。平衡でどこまで溶けるかと、そこへ至る速度、風化抵抗という経験系列を別の概念として整理する。

\Needspace{4\baselineskip}\paragraph{選択肢の検討（査読待ち初稿）}\mbox{}\par

\noindent\textbf{1：不適切} Debye-Huckelは造岩鉱物の風化抵抗推定モデル。。Debye-Huckelは希薄溶液の静電相互作用と活量係数を扱う。風化抵抗系列そのもののモデルではない。\par

\noindent\textbf{2：適切} 各イオン濃度からイオン強度を計算。。価数zと濃度cを用いて\(I=\frac12\sum_i c_i z_i^2\)。価数の二乗を落とさず濃度基準を揃える。\par

\noindent\textbf{3：成立条件に疑義} 溶解の最終平衡濃度が溶解度。。固定温度・圧力・化学条件で残存固相と平衡に達した飽和濃度である。固相が全部溶けた場合や流動開放系はその濃度が溶解度とは限らない。\par

\noindent\textbf{4：適切（経験系列）} 風化抵抗はかんらん石→輝石→黒雲母→石英の順。。地表で高温形成鉱物が相対的に不安定な通常の風化系列の一部である。粒径や欠陥、流体条件で実際の速度は変わる。\par

\noindent\textbf{5：適切（通常反応条件）} 方解石はケイ酸塩鉱物より溶解が速い。。通常の炭酸を含む地表水では方解石は比較的速く溶ける。平衡に近いと速度は小さく、全温度・pH・鉱物に普遍とはしない。\par

\paragraph{再利用する考え方} 物質の貯留量・輸送速度・化学形態を分け、系の境界と反応の前提を確かめる。溶解の平衡・活量と反応速度の分離を収支・平衡・反応速度のどのモデルで扱うか選ぶ。

\paragraph{共通理論} \hyperref[theory:geochemistry:activity]{第\ref*{theory:geochemistry:activity}節}, \hyperref[theory:geochemistry:equilibrium]{第\ref*{theory:geochemistry:equilibrium}節}, \hyperref[theory:geochemistry:weathering]{第\ref*{theory:geochemistry:weathering}節}

\par 参考：\cite{gc_debye,gc_hem,gc_clays}。

\paragraph{関連問題} \hyperref[q:r06-reg-applied-22]{R06-REG-APPLIED-22}

\par 難易度 3（校準初稿）、優先度未評価。

\subsection{令和6年度 問22}
\label{q:r06-reg-applied-22}

\noindent\textbf{原典：} sources/R06/17\_applied\_science.pdf、PDF物理7頁。

\textbf{公式正答：} 2（single）。

\paragraph{問題の要旨} 土壌の酸化還元、炭酸塩析出、乾燥塩類、金属沈殿と粘土生成を比較する。

\Needspace{4\baselineskip}\paragraph{確認した条件・記述（独自要約）}\mbox{}\par

\noindent\textbf{方向：} 最も不適切なもの。\par

\noindent\textbf{対象：} 地球表層の自然水・岩石・元素循環。各肢の条件は独自要約として保持。\par

\paragraph{問われる核心} 化学的風化・炭酸平衡と土壌の二次鉱物

\paragraph{必要知識} 炭酸と炭酸水素イオンの平衡、CO2脱気による方解石析出、毛管上昇と蒸発濃縮、Al・Feの酸化状態と溶解度、溶脱と二次粘土鉱物。

\paragraph{考え方・正答を選ぶ理由（査読待ち初稿）} 2のCO2添加が逆である。\[\mathrm{CaCO_3+CO_2+H_2O\rightleftharpoons Ca^{2+}+2HCO_3^-}\]においてCO2を失うと左へ進み方解石を析出する。溶解、沈殿、蒸発濃縮、酸化還元、強溶脱の作用をそれぞれ区別して他肢を通常条件で判断する。

\Needspace{4\baselineskip}\paragraph{選択肢の検討（査読待ち初稿）}\mbox{}\par

\noindent\textbf{1：適切（通常条件）} 土壌有機物は還元型。。CO2に比べ有機炭素は還元されており、酸化で電子を供給できる。\par

\noindent\textbf{2：不適切} 石灰岩溶解地下水へ大気CO2が加わって鍾乳石・石筍形成。。普通の石灰岩洞窟ではCO2の脱気でCaCO3が析出する。添加では炭酸が増して溶解側へ動くので方向が逆。\par

\noindent\textbf{3：適切（通常条件）} 乾燥地で塩類が上昇・蒸発沈積。。毛管上昇と蒸発で溶存塩が濃縮される。原材料と飽和条件に応じ岩塩・石膏・炭酸塩等が析出する。\par

\noindent\textbf{4：適切（通常条件）} 通常土壌pHでAl3+やFeは酸化物・オキシ水酸化物沈殿。。通常の酸化的・弱酸性〜中性土壌ではAl・Fe(III)の水酸化物等が溶解を制限する。強酸、錯体、還元Fe(II)などの例外を付す。\par

\noindent\textbf{5：適切（通常条件）} 湿潤熱帯の強風化ではスメクタイトよりカオリナイト。。強い溶脱は塩基性陽イオンと一部Siを除き、カオリナイト等の比較的単純な相を残しやすい。気候だけで鉱物を一意に決めない。\par

\paragraph{再利用する考え方} 物質の貯留量・輸送速度・化学形態を分け、系の境界と反応の前提を確かめる。化学的風化・炭酸平衡と土壌の二次鉱物を収支・平衡・反応速度のどのモデルで扱うか選ぶ。

\paragraph{共通理論} \hyperref[theory:geochemistry:carbonate]{第\ref*{theory:geochemistry:carbonate}節}, \hyperref[theory:geochemistry:weathering]{第\ref*{theory:geochemistry:weathering}節}, \hyperref[theory:geochemistry:redox]{第\ref*{theory:geochemistry:redox}節}

\par 参考：\cite{gc_cave,gc_clays,gc_hem,gc_redox}。

\paragraph{関連問題} \hyperref[q:r04-reg-applied-20]{R04-REG-APPLIED-20}

\par 難易度 3（校準初稿）、優先度未評価。

\section{水質汚濁指標と起源トレーサー}

\subsection{平成23年度 問18}
\label{q:h23-reg-applied-18}

\noindent\textbf{原典：} sources/H23/17\_applied\_science.pdf、PDF物理6頁。

\textbf{公式正答：} 5（single）。

\paragraph{問題の要旨} 汚濁量を表す水質指標と、水の起源を追う同位体比を区別する。

\Needspace{4\baselineskip}\paragraph{確認した条件・記述（独自要約）}\mbox{}\par

\noindent\textbf{方向：} 最も不適切なもの。\par

\noindent\textbf{対象：} 地球表層の自然水・岩石・元素循環。各肢の条件は独自要約として保持。\par

\paragraph{問われる核心} 水質指標と起源トレーサーの用途

\paragraph{必要知識} 水質指標と起源トレーサーの用途。汚濁を評価する量と水の由来を同定する量は用途が異なる。1〜4は衛生・有機負荷・栄養塩・粒子の水質指標で、5が最も不適切である。同位体が汚染研究に利用できないという意味ではなく、単独の比を汚濁負荷へ変換できないという区別である。

\paragraph{考え方・正答を選ぶ理由（査読待ち初稿）} 汚濁を評価する量と水の由来を同定する量は用途が異なる。1〜4は衛生・有機負荷・栄養塩・粒子の水質指標で、5が最も不適切である。同位体が汚染研究に利用できないという意味ではなく、単独の比を汚濁負荷へ変換できないという区別である。

\Needspace{4\baselineskip}\paragraph{選択肢の検討（査読待ち初稿）}\mbox{}\par

\noindent\textbf{1：適切（通常条件）} 大腸菌群数。当時の衛生汚濁の指標である。菌数を物質濃度や発生源の同位体指標と混ぜない。\par

\noindent\textbf{2：適切（通常条件）} 生物化学的酸素要求量BOD。微生物分解に必要な酸素量で、分解可能な有機負荷を反映する。全有機物量そのものではない。\par

\noindent\textbf{3：適切（通常条件）} 全窒素。富栄養化に関与する窒素成分の総量を追う。\par

\noindent\textbf{4：適切（通常条件）} 浮遊物質SS。水に浮遊する粒子量を追う。溶存量とは区別する。\par

\noindent\textbf{5：不適切} 水素同位体比。蒸発・降水・混合等による水の起源と履歴を追うトレーサーであり、一般的な汚濁量の直接尺度ではない。\par

\paragraph{再利用する考え方} 物質の貯留量・輸送速度・化学形態を分け、系の境界と反応の前提を確かめる。水質指標と起源トレーサーの用途を収支・平衡・反応速度のどのモデルで扱うか選ぶ。

\paragraph{共通理論} \hyperref[theory:geochemistry:indicators]{第\ref*{theory:geochemistry:indicators}節}, \hyperref[theory:geochemistry:isotope]{第\ref*{theory:geochemistry:isotope}節}

\par 参考：\cite{gc_hem,gc_meteoric}。

\par 難易度 3（校準初稿）、優先度未評価。

\subsection{平成30年度 問28}
\label{q:h30-reg-applied-28}

\noindent\textbf{原典：} sources/H30/17\_applied\_science.pdf、PDF物理9頁。

\textbf{公式正答：} 2（single）。

\par\textbf{検討未確定：} 肢2の単位誤りは確認。肢3「金属ヒ素に毒性はない」は無条件に支持できず、肢4の岩種量比較と揮発性の因果は一次定量未確認。2018原典を科学的に正しい文へ修正せず疑義保持。

\paragraph{問題の要旨} 自然ヒ素、水道基準の単位、単体の有害性と岩石・鉱物での存在を比較する。

\Needspace{4\baselineskip}\paragraph{確認した条件・記述（独自要約）}\mbox{}\par

\noindent\textbf{方向：} 最も不適切。\par

\noindent\textbf{時点：} 2018年出題。2004施行表・2026年閲覧現行表を区別。\par

\noindent\textbf{単位：} 原典肢2はg/Lでありmg/Lへ校正しない。\par

\paragraph{問われる核心} 自然ヒ素の形態と濃度単位・基準の時点

\paragraph{必要知識} 自然ヒ素の形態と濃度単位・基準の時点。公式2・初回2を保持する。水道基準0.01mg/Lに対し原典は0.01g/Lであり1000倍違うため、2の不適切さは独立に確認できる。ただし3の無毒という断言は現在の一次有害性資料で支持できず、4の岩種比較も未確定である。正答だけが合ったことを全肢の科学的妥当性へ昇格させない。

\paragraph{考え方・正答を選ぶ理由（査読待ち初稿）} 公式2・初回2を保持する。水道基準0.01mg/Lに対し原典は0.01g/Lであり1000倍違うため、2の不適切さは独立に確認できる。ただし3の無毒という断言は現在の一次有害性資料で支持できず、4の岩種比較も未確定である。正答だけが合ったことを全肢の科学的妥当性へ昇格させない。

\Needspace{4\baselineskip}\paragraph{選択肢の検討（査読待ち初稿）}\mbox{}\par

\noindent\textbf{1：適切} 鉱物由来無機ヒ素が温泉・地下水へ溶出。。鉱物反応や酸化還元・吸着の変化によって自然由来の無機ヒ素が溶出する場合がある。\par

\noindent\textbf{2：不適切（単位1000倍）} 水道基準は0.01g/L以下。。2004施行表と現行表ともヒ素量0.01mg/Lで、0.01g/Lは10mg/L、1000倍の単位誤りである。\par

\noindent\textbf{3：無毒の断言に疑義} 半導体用金属ヒ素に毒性はない。。単体は溶解・取り込みが化合物と異なるが、NIOSHは金属ヒ素も有害性と曝露制限の対象として示す。無毒という絶対命題を教材の安全根拠にしない。\par

\noindent\textbf{4：比較と因果は要確認} 揮発性のため火成・変成岩にはあまり含まれない。。ヒ素は岩石中に微量含まれ火山ガスにも存在する。岩種ごとの含有量、鉱化、揮発移動を分ける必要があり、原典の因果・比較範囲は追加確認が必要である。\par

\noindent\textbf{5：適切（形態による）} 地表の水環境で亜ヒ酸塩・ヒ酸塩鉱物が形成。。酸化数III・Vに応じた亜ヒ酸塩・ヒ酸塩等の形態と二次鉱物があり、水と鉱物の反応に関わる。溶存種と固相鉱物を混ぜない。\par

\paragraph{再利用する考え方} 物質の貯留量・輸送速度・化学形態を分け、系の境界と反応の前提を確かめる。自然ヒ素の形態と濃度単位・基準の時点を収支・平衡・反応速度のどのモデルで扱うか選ぶ。

\paragraph{共通理論} \hyperref[theory:geochemistry:arsenic]{第\ref*{theory:geochemistry:arsenic}節}, \hyperref[theory:geochemistry:redox]{第\ref*{theory:geochemistry:redox}節}

\par 参考：\cite{gc_hem,gc_arsenic2004,gc_arsenic_current,gc_arsenic_hazard}。

\par 難易度 3（校準初稿）、優先度未評価。

